% =====================================================================
\chapter{Existencia, unicidad, contraejemplos y redacción}
\label{tema:6}
% =====================================================================

Los temas anteriores se han centrado en enunciados universales. Este tema trata los enunciados existenciales y de unicidad, explica cómo se refuta un enunciado falso y termina con unas pautas para redactar demostraciones claras, que resumen todo el módulo. Se siguen los capítulos~7 y~9 y la §5.3 de~\cite{hammack}, la §3.6 de~\cite{velleman}, la §1.9 de~\cite{mcs}, la lista de comprobación de CS103~\cite{cs103} y el método de Polya~\cite{polya}.

\section{Demostraciones de existencia}

Para demostrar $\exists x\,P(x)$ basta encontrar \emph{un} elemento $a$, llamado \emph{testigo}, y comprobar $P(a)$. Una demostración de este tipo es \emph{constructiva}. A veces se prueba la existencia sin exhibir el testigo, por ejemplo por reducción al absurdo o examinando casos sin saber cuál se da; la demostración es entonces \emph{no constructiva}~\cite[§7.3–7.4]{hammack}.

\begin{ejemplo}[Existencia constructiva]\label{ej:1729}
Existe un entero positivo que es suma de dos cubos positivos de dos maneras distintas: $1729=1^{3}+12^{3}=9^{3}+10^{3}$. La demostración consiste en dar el número y hacer las cuentas: $1+1728=1729$ y $729+1000=1729$.
\end{ejemplo}

\begin{proposicion}[Existencia no constructiva]\label{prop:irracional-potencia}
Existen números irracionales $a$ y $b$ tales que $a^{b}$ es racional.
\end{proposicion}

\begin{proof}
Sea $c=\sqrt2^{\sqrt2}$. Se distinguen dos casos. Si $c$ es racional, basta tomar $a=b=\sqrt2$, irracionales por el \cref{teo:raiz2}. Si $c$ es irracional, se toman $a=c$ y $b=\sqrt2$, y entonces
\[
a^{b}=\bigl(\sqrt2^{\sqrt2}\bigr)^{\sqrt2}=\sqrt2^{\sqrt2\cdot\sqrt2}=\sqrt2^{\,2}=2,
\]
que es racional. En ambos casos existen $a$ y $b$ con la propiedad.
\end{proof}

La prueba no dice cuál de los dos pares sirve. (Se sabe, por un teorema profundo de Gelfond y Schneider, que $c$ es irracional, pero la demostración anterior no lo necesita.)

\section{Demostraciones de unicidad}

Para demostrar $\exists!\,x\,P(x)$ se prueban dos cosas (\cref{def:unicidad}): la \emph{existencia}, $\exists x\,P(x)$, y la \emph{unicidad}, que se redacta así: «sean $x$ e $y$ tales que $P(x)$ y $P(y)$; entonces $x=y$». Una alternativa equivalente es suponer $x\neq y$ y llegar a una contradicción~\cite[§3.6]{velleman}.

\begin{proposicion}[Unicidad en la división euclídea]\label{prop:division-unicidad}
Sean $a\in\Z$ y $b\in\N$. Si $a=qb+r=q'b+r'$ con $q,q',r,r'\in\Z$ y $0\le r,r'<b$, entonces $q=q'$ y $r=r'$.
\end{proposicion}

\begin{proof}
Restando las dos expresiones, $(q-q')b=r'-r$. De $0\le r,r'<b$ se sigue $-b<r'-r<b$, es decir, $\abs{q-q'}\,b<b$, luego $\abs{q-q'}<1$. Como $q-q'$ es entero, $q=q'$, y entonces $r'-r=0$.
\end{proof}

Junto con la \cref{prop:division-existencia}, esto completa la demostración del \cref{teo:division}.

\begin{ejemplo}\label{ej:unico-real}
La ecuación $x^{3}+x-2=0$ tiene una única solución real. \emph{Existencia:} $x=1$ cumple $1+1-2=0$. \emph{Unicidad:} $x^{3}+x-2=(x-1)(x^{2}+x+2)$, y el segundo factor no se anula en $\R$ porque su discriminante es $1-8=-7<0$. Si $x^{3}+x-2=0$, el producto se anula y, como el segundo factor no es cero, $x-1=0$.
\end{ejemplo}

\section{Contraejemplos y refutaciones}

Para \emph{refutar} un enunciado se demuestra su negación~\cite[cap.~9]{hammack}. Según la forma del enunciado:
\begin{itemize}[leftmargin=1.4em]
  \item $\forall x\,P(x)$ se refuta con un \emph{contraejemplo}: un $a$ concreto con $\neg P(a)$.
  \item $\forall x\,(P(x)\Rightarrow Q(x))$ se refuta con un $a$ que cumple $P(a)$ y no cumple $Q(a)$.
  \item $\exists x\,P(x)$ se refuta demostrando $\forall x\,\neg P(x)$, lo que exige una demostración general, no un ejemplo.
\end{itemize}

\begin{ejemplo}[El polinomio de Euler]\label{ej:euler}
El polinomio $f(n)=n^{2}+n+41$ toma valores primos para $n=0,1,\dots,39$, lo que podría sugerir que es primo para todo $n\in\N_0$. El enunciado es falso: $f(40)=1600+40+41=1681=41^{2}$. Muchas comprobaciones favorables no demuestran un enunciado universal; un único contraejemplo lo refuta.
\end{ejemplo}

\begin{ejemplo}\label{ej:divide-producto}
«Si $a\divides bc$, entonces $a\divides b$ o $a\divides c$» es falso para enteros cualesquiera: con $a=6$, $b=4$, $c=9$ se tiene $6\divides 36$, pero $6\nodivide 4$ y $6\nodivide 9$. (Es cierto cuando $a$ es primo; es el lema de Euclides del \cref{eje:4-raiz-primo}.)
\end{ejemplo}

\section{Cómo redactar una demostración}

Una demostración es un texto dirigido a un lector, no una lista de cálculos. Las pautas siguientes resumen las de la §5.3 de~\cite{hammack}, la §1.9 de~\cite{mcs} y la lista de comprobación de CS103~\cite{cs103}.

\begin{enumerate}[label=\arabic*.]
  \item \emph{Anunciar el plan.} Indicar qué se supone, qué se quiere demostrar y con qué técnica («Se procede por contrarrecíproco», «Por inducción sobre $n$»).
  \item \emph{Introducir cada variable.} Toda letra debe presentarse antes de usarse, con su universo: «Sea $n$ un entero impar»; «existe $k\in\Z$ con\dots». No reutilizar una letra para objetos distintos.
  \item \emph{Usar las definiciones, no repetirlas.} Desplegar la definición en el punto donde se necesita, aplicada a los objetos concretos.
  \item \emph{Escribir frases completas.} Las fórmulas forman parte de las frases. Los símbolos $\forall$, $\exists$, $\Rightarrow$ son para enunciados formales, no abreviaturas en la prosa.
  \item \emph{Que cada frase aporte.} Eliminar lo que no se usa después; justificar cada paso no evidente citando el resultado empleado.
  \item \emph{Terminar con la conclusión.} Decir explícitamente que se ha demostrado lo que se buscaba y cerrar con $\blacksquare$.
\end{enumerate}

Para \emph{encontrar} la demostración, antes de redactarla, sirve el método de cuatro fases de Polya~\cite{polya}: comprender el problema (datos, incógnita, condiciones; escribir el enunciado con cuantificadores), concebir un plan (¿se conoce un problema relacionado?, ¿sirve un caso particular o el contrarrecíproco?), ejecutarlo comprobando cada paso, y examinar la solución obtenida (¿se usaron todas las hipótesis?, ¿hay un argumento más directo?).

\begin{ejemplo}[Borrador y redacción final]\label{ej:redaccion}
Enunciado: «si $a$ y $b$ son enteros impares, $a+b$ es par». Un borrador típico es «$a=2k+1$, $b=2k+1$, $a+b=4k+2=2(2k+1)$ par». Tiene dos defectos: usa la misma $k$ para dos números distintos y no presenta las variables. Redacción correcta: «Sean $a$ y $b$ enteros impares. Por definición existen $k,\ell\in\Z$ con $a=2k+1$ y $b=2\ell+1$. Entonces $a+b=2k+2\ell+2=2(k+\ell+1)$, y como $k+\ell+1$ es entero, $a+b$ es par. $\blacksquare$»
\end{ejemplo}

\begin{observacion}[Errores frecuentes]\label{obs:6-errores}
(a) «Demostrar» un enunciado universal con ejemplos (\cref{ej:euler}).
(b) Refutar un enunciado existencial con un ejemplo en que falla: que $P(5)$ sea falsa no impide que exista otro $x$ con $P(x)$.
(c) En una prueba de unicidad, suponer de entrada que los dos elementos son distintos y olvidar después usar esa suposición.
(d) Escribir una cadena de igualdades sin indicar cuál es el punto de partida y cuál el de llegada.
\end{observacion}

\begin{observacion}[Conexión con el curso]
Los resultados de existencia y unicidad son frecuentes en el curso: el sistema de ecuaciones normales de mínimos cuadrados tiene solución única cuando las columnas de la matriz de diseño son linealmente independientes (módulos~2 y~8), y una función estrictamente convexa tiene a lo sumo un mínimo (módulo~7). Los contraejemplos aparecen para marcar los límites de un resultado: el mínimo de dos funciones convexas no es convexo (módulo~7), el criterio de la hessiana no decide si hay un autovalor nulo (módulo~4) y la fórmula de la longitud tras una fusión de BPE falla cuando el par se solapa consigo mismo (módulo~12).
\end{observacion}

\begin{lecturas}
\item Hammack~\cite{hammack}, cap.~7 «Proving Non-Conditional Statements» (§7.3–7.4), cap.~9 «Disproof» y §5.3 «Mathematical Writing».
\item Velleman~\cite{velleman}, §3.6 «Existence and Uniqueness Proofs» y §3.7 «More Examples of Proofs».
\item Lehman, Leighton y Meyer~\cite{mcs}, §1.9 «Good Proofs in Practice».
\item Stanford CS103~\cite{cs103}, «Proofwriting Checklist».
\item Polya~\cite{polya}, primera parte (las cuatro fases de la resolución de problemas).
\end{lecturas}

\section*{Ejercicios del tema 6}
\addcontentsline{toc}{section}{Ejercicios del tema 6}

\begin{ejercicio}{1}{eje:6-clasificar}
Para cada enunciado, indique si es universal, existencial o de existencia y unicidad, y qué haría falta para demostrarlo y para refutarlo: (a) «todo primo mayor que~2 es impar»; (b) «hay un primo entre $100$ y $110$»; (c) «la ecuación $2x+3=7$ tiene una única solución real»; (d) «ningún entero cumple $n^{2}=3$».
\end{ejercicio}

\begin{ejercicio}{2}{eje:6-contraejemplos}
Refute con un contraejemplo: (a) para todo $x\in\R$, $x^{2}>x$; (b) para todos $a,b\in\R$, $\sqrt{a^{2}+b^{2}}=a+b$; (c) si $n\in\N$ y $n^{2}$ es múltiplo de~4, entonces $n$ es múltiplo de~4; (d) para todo $n\in\N$, $2^{n}+1$ es primo.
\fuente{adaptado de Hammack~\cite{hammack}, §9.1 «Counterexamples».}
\end{ejercicio}

\begin{ejercicio}{2}{eje:6-demostrar-refutar}
Demuestre o refute: (a) la suma de dos números irracionales es irracional; (b) el producto de dos números racionales es racional; (c) existe $n\in\N$ tal que $n^{2}+n+1$ es par; (d) para todo $x\in\R$ existe $y\in\R$ con $y^{2}=x$.
\fuente{adaptado de Hammack~\cite{hammack}, cap.~9 «Disproof» (§9.1–9.2).}
\end{ejercicio}

\begin{ejercicio}{3}{eje:6-testigos}
Demuestre: (a) existen números reales $a$ y $b$ con $\floor{a}\cdot\ceil{b}\neq\floor{ab}$; (b) existen naturales $a>b>c>d>0$ con $a^{2}+b^{2}+c^{2}+d^{2}=137$.
\fuente{adaptado de Stanford CS103~\cite{cs103}, \emph{Problem Set}~1, problema~7.}
\end{ejercicio}

\begin{ejercicio}{3}{eje:6-irracional-entre}
Sean $r<s$ números racionales. Demuestre que existe un número irracional $x$ con $r<x<s$.
\end{ejercicio}

\begin{ejercicio}{3}{eje:6-paridad-unica}
Demuestre que para cada $n\in\Z$ existe un único $k\in\Z$ tal que $n=2k$ o $n=2k+1$, y que ningún entero es a la vez par e impar. (Puede usar el \cref{teo:division}; para la segunda parte dé además una demostración directa por reducción al absurdo.)
\end{ejercicio}

\begin{ejercicio}{3}{eje:6-terna}
Demuestre que existe un único número primo $p$ tal que $p$, $p+2$ y $p+4$ son todos primos.
\end{ejercicio}

\begin{ejercicio}{3}{eje:6-mersenne}
Demuestre o refute: (a) si $n$ es primo, entonces $2^{n}-1$ es primo; (b) si $2^{n}-1$ es primo, entonces $n$ es primo (con $n\ge 2$). (Indicación para (b): por contrarrecíproco, use que $x^{m}-1=(x-1)(x^{m-1}+\dots+x+1)$.)
\end{ejercicio}

\begin{ejercicio}{2}{eje:6-promedio}
Sean $x_1,\dots,x_{10}$ números reales con $x_1+\dots+x_{10}=100$. Demuestre que existe $i$ con $x_i\ge 10$. ¿Es constructiva la demostración? ¿Existe necesariamente $i$ con $x_i>10$?
\end{ejercicio}

\begin{ejercicio}{2}{eje:6-critica-1}
Analice esta prueba de «la suma de un entero par y uno impar es impar»: «Sea $x=2k$ par e $y=2k+1$ impar. Entonces $x+y=4k+1$. Por tanto, $x+y$ es impar.» Identifique los defectos de redacción y de razonamiento, decida si el teorema es cierto y redacte una demostración correcta.
\fuente{adaptado de Stanford CS103~\cite{cs103}, \emph{Problem Set}~1, problema~9 «Proof Critiques».}
\end{ejercicio}

\begin{ejercicio}{2}{eje:6-critica-2}
Analice esta prueba de «para todo $n\in\Z$, si $n\equiv 0\pmod 2$, entonces $n\equiv 0\pmod 4$»: «Se procede por contrarrecíproco. Supóngase $n\equiv 0\pmod 4$. Entonces $n=4k=2(2k)$, luego $n\equiv 0\pmod 2$.» Indique qué se ha demostrado en realidad y si el enunciado es cierto.
\fuente{adaptado de Stanford CS103~\cite{cs103}, \emph{Problem Set}~1, problema~9 «Proof Critiques».}
\end{ejercicio}

\begin{ejercicio}{2}{eje:6-uno-dos}
Encuentre el error: «Sean $a=b$. Entonces $a^{2}=ab$, luego $a^{2}-b^{2}=ab-b^{2}$, es decir, $(a+b)(a-b)=b(a-b)$. Simplificando, $a+b=b$; como $a=b$, $2b=b$ y, dividiendo por $b$, $2=1$.»
\end{ejercicio}

\begin{ejercicio}{3}{eje:6-redactar}
El siguiente texto contiene las ideas de una demostración, pero no está redactado como tal: «$a\divides b$, $b\divides a$, $a,b>0$ $\Rightarrow$ $b=ak$, $a=bl$ $\Rightarrow$ $a=akl$ $\Rightarrow$ $kl=1$ $\Rightarrow$ $k=l=1$ $\Rightarrow$ $a=b$.» Redáctelo como una demostración completa de «si $a,b\in\N$, $a\divides b$ y $b\divides a$, entonces $a=b$», justificando especialmente el paso $kl=1\Rightarrow k=l=1$.
\fuente{adaptado de Hammack~\cite{hammack}, §5.3 «Mathematical Writing».}
\end{ejercicio}

\begin{ejercicio}{3}{eje:6-unica-raiz}
Sean $a,b,c\in\R$ con $a\neq 0$. Demuestre que la ecuación $ax^{2}+bx+c=0$ tiene una única solución real si y solo si $b^{2}=4ac$, y que en ese caso la solución es $x=-b/(2a)$.
\end{ejercicio}
